% =====================================================================
%  STANDALONE PAPER.  Compile with:  pdflatex hierarchy_paper.tex  (x2)
%  The body is identical to hierarchy_body.tex, which can instead be
%  \input into the existing Phase 3 write-up.
% =====================================================================
\documentclass[11pt,a4paper]{article}
\usepackage[margin=1in]{geometry}
\usepackage{amsmath,amssymb,amsthm}
\usepackage{booktabs}
\usepackage[hidelinks]{hyperref}
%\usepackage{microtype}  % disabled: needs scalable fonts

\theoremstyle{plain}
\newtheorem{theorem}{Theorem}
\newtheorem{proposition}[theorem]{Proposition}
\theoremstyle{definition}
\newtheorem{definition}[theorem]{Definition}
\newtheorem{construction}[theorem]{Construction}

\title{Recovering latent hierarchy without gradient descent:\\
a negative result for transformers, a verified estimator,\\
and calibrated instruments for detection}
\author{}
\date{}

\begin{document}
\maketitle

\begin{abstract}
We study recovery of the latent symbols of a depth-$4$ probabilistic
context-free grammar from its leaf strings. Three results. First, across twelve
transformer configurations---varying depth, width, input representation, output
factorisation, architecture family, auxiliary supervision and
curriculum---a linear probe for the level-3 latent symbol never leaves its
shuffled-label floor, although an exact computation shows $1.87$ nats of
level-3 information are present in the observed prefix. Second, a gradient-free
estimator recovers the grammar to $0.001$ nats per sequence of the true model
in $131$ seconds, once two things are fixed: the regime must be identifiable
(we give the observations-per-parameter arithmetic separating $0.15$ from
$781$), and the clustering must happen \emph{inside} EM rather than being
imposed beforehand. A hard partition of observed child pairs cannot represent a
grammar in which pairs are shared across parents, which $4$ of $16$ level-1
pairs are here; split-merge moves scored by Bregman information close the
remaining gap. Third, we build and validate instruments for detecting
hierarchy and its failure: a parametric-bootstrap obstruction test with a
planted positive ($z=+220$ against $+6.5$ elsewhere), a block-level dependence
measure, and a $\beta_1$ test for non-tree structure that fires on a
local-optimum fit ($6$ of $25$ thresholds), is silent on the recovered one
($0$ of $25$), and catches planted couplings in both cases ($9$ of $25$).
Every structural identity used is verified numerically to $10^{-14}$. We also
record what does not work, including a probe result we withdraw, the vacuity of
the cohomological gluing obstruction on a single tree, and a proof that the
proposed $p$-adic lattice quasi-morphism is constant.
\end{abstract}

\tableofcontents
\bigskip

% =====================================================================
%  Hierarchy recovery: transformer failure, constructed estimator,
%  and verified detection instruments.
%
%  INSERTABLE BODY.  Drop into an existing document with
%      \input{hierarchy_body}
%  Requires in the host preamble: amsmath, amssymb, booktabs, hyperref.
%  Uses \section as top level; change to \subsection if inserting under
%  an existing Phase 3 section.
% =====================================================================

% =====================================================================
%  Hierarchy recovery: transformer failure, constructed estimator,
%  and verified detection instruments.
%
%  INSERTABLE BODY.  Drop into an existing document with
%      \input{hierarchy_body}
%  Requires in the host preamble: amsmath, amssymb, booktabs, hyperref.
%  Uses \section as top level; change to \subsection if inserting under
%  an existing Phase 3 section.
% =====================================================================

% =====================================================================
%  Hierarchy recovery: transformer failure, constructed estimator,
%  and verified detection instruments.
%
%  INSERTABLE BODY.  Drop into an existing document with
%      \input{hierarchy_body}
%  Requires in the host preamble: amsmath, amssymb, booktabs, hyperref.
%  Uses \section as top level; change to \subsection if inserting under
%  an existing Phase 3 section.
% =====================================================================

% =====================================================================
%  Hierarchy recovery: transformer failure, constructed estimator,
%  and verified detection instruments.
%
%  INSERTABLE BODY.  Drop into an existing document with
%      \input{hierarchy_body}
%  Requires in the host preamble: amsmath, amssymb, booktabs, hyperref.
%  Uses \section as top level; change to \subsection if inserting under
%  an existing Phase 3 section.
% =====================================================================

% =====================================================================
%  Hierarchy recovery: transformer failure, constructed estimator,
%  and verified detection instruments.
%
%  INSERTABLE BODY.  Drop into an existing document with
%      \input{hierarchy_body}
%  Requires in the host preamble: amsmath, amssymb, booktabs, hyperref.
%  Uses \section as top level; change to \subsection if inserting under
%  an existing Phase 3 section.
% =====================================================================

% =====================================================================
%  Hierarchy recovery: transformer failure, constructed estimator,
%  and verified detection instruments.
%
%  INSERTABLE BODY.  Drop into an existing document with
%      \input{hierarchy_body}
%  Requires in the host preamble: amsmath, amssymb, booktabs, hyperref.
%  Uses \section as top level; change to \subsection if inserting under
%  an existing Phase 3 section.
% =====================================================================

\section{Setup: the generative model and its exact floors}
\label{sec:hier-setup}

All experiments use the Random Hierarchy Model (RHM) of
Cagnetta et al.~\cite{cagnetta2024}: a probabilistic context-free grammar on a
fixed complete binary tree of depth $L$, with $V$ symbols per internal level,
$m$ productions per symbol, and a leaf alphabet of size $A$. A sequence is the
string of $2^L$ leaves obtained by sampling a root symbol uniformly and then
one production per node. Because the tree shape is fixed and positional, the
latent structure is a Markov chain
\begin{equation}
  \text{leaves} \;\to\; \text{level }1 \;\to\; \cdots \;\to\; \text{root},
  \label{eq:markov}
\end{equation}
and every conditional independence we use below is an instance of the tree
Markov property.

Two corpora appear. Corpus~A ($L=4$, $V=64$, $A=48$, $m=12$, $2\times10^5$
sequences of length $16$) is the configuration inherited from earlier phases.
Corpus~B ($L=4$, $V=8$, $A=12$, $m=3$, $2\times10^5$ sequences) was built for
this study; Section~\ref{sec:hier-regime} explains why.

Positions are labelled \emph{local} (odd indices, determined by the sibling)
and \emph{higher} (even indices, which open a new subtree). We compute the
\emph{exact} per-position conditional entropy $H(x_t \mid x_{<t})$ from the
grammar by the inside--outside algorithm (Theorem~\ref{thm:inside}), rather
than estimating it. For Corpus~A the means are
\begin{equation}
  H_{\mathrm{local}} = 1.6034 \ \text{nats}, \qquad
  H_{\mathrm{higher}} = 3.5527 \ \text{nats},
  \label{eq:floorsA}
\end{equation}
with $H(x_0)=3.8161$ alone (the first token of a fresh tree is nearly
unpredictable), and $3.5150$ for the higher positions excluding $t=0$. These
are floors: no predictor of any kind can go below them.

We also compute, exactly, the conditional entropy of each \emph{latent} symbol
given the observed prefix. Over the higher positions of Corpus~A, against a
chance value of $\ln V = 4.1589$:

\begin{center}
\begin{tabular}{lcccc}
\toprule
level & 1 & 2 & 3 & 4 \\
\midrule
$H(\text{symbol} \mid x_{<t})$ & 2.7835 & 2.2923 & 2.2900 & 2.7910 \\
fraction of chance & 67\% & 55\% & 55\% & 67\% \\
\bottomrule
\end{tabular}
\end{center}

\noindent
At individual positions the level-3 symbol is nearly determined ($0.427$ nats
at $t=14$, $1.126$ at $t=12$). This matters for what follows: roughly $1.87$
nats of level-3 information are recoverable from the prefix, so a model that
fails to represent it is failing at something the data support, not attempting
the impossible.


\section{Part I: transformers do not recover levels 3--4}
\label{sec:hier-transformer}

\subsection{Protocol}

Decoder-only transformers were trained on Corpus~A with next-token
cross-entropy for $48{,}000$ steps. After training, a linear probe was fitted
from the residual stream $h[t]$ to the level-$\ell$ symbol above position $t$,
with a $70/30$ split over sequences and two controls: shuffled labels, and an
untrained network of identical shape. Chance is $1/V = 0.0156$.

\subsection{Result}

\begin{center}
\begin{tabular}{lccc}
\toprule
configuration & \texttt{local} & \texttt{higher} & level-3 probe \\
\midrule
2 layers, $d=64$ (baseline)      & 2.173 & 3.831 & 0.028 \\
4 layers                         & 2.078 & 3.797 & --- \\
8 layers, frozen embedding       & 2.037 & 3.755 & --- \\
8 layers, frozen emb., fresh sampling & 2.022 & 3.772 & 0.027 \\
$d=128$, 4 layers                & 2.039 & 3.770 & --- \\
frozen random embedding (2 layers) & 2.186 & 3.824 & --- \\
constructed order coordinate     & 2.162 & 3.815 & --- \\
auxiliary supervision, levels 3--4 & 2.088 & 3.777 & 0.056 \\
latent marginalising head (stopped at 34k) & 3.227 & 3.835 & --- \\
dilated causal CNN (tree receptive fields) & 2.104 & 3.827 & 0.032 \\
curriculum over tree depth       & 2.085 & 3.803 & 0.034 \\
curriculum control (no easy cases) & 2.043 & 3.789 & 0.033 \\
\midrule
exact floor                      & \textbf{1.603} & \textbf{3.553} & --- \\
\bottomrule
\end{tabular}
\end{center}

\noindent
Shuffled-label controls ran at $0.021$--$0.034$; the level-3 probe is
indistinguishable from them in every configuration. Level~4 read $0.016$--$0.020$
against chance $0.0156$. The auxiliary arm is the strongest test: its head was
trained jointly for $48{,}000$ steps on \emph{ground-truth} level-3 and level-4
labels and reached $0.056$ and $0.019$.

Three additional controls sharpen the claim.

\paragraph{The embedding is not the site.}
With the token embedding frozen at random initialisation and never updated,
final loss is unchanged ($2.185/3.824$ versus $2.163/3.831$), and the
sublayer readout shows the order structure being constructed entirely in the
blocks: the embedding tap reads exactly $0.0000$, then $+0.0430$ (attn$_0$),
$+0.0508$ (mlp$_0$), $+0.0410$ (attn$_1$), $+0.1272$ (mlp$_1$).

\paragraph{A linear probe on the embedding was not evidence.}
An earlier phase reported a linear probe reading comparability off the
embedding matrix at $0.98$--$1.00$ on held-out multi-hop pairs. Running the
identical probe on (i) a fresh random initialisation, (ii) an exactly
orthonormal matrix, and (iii) the trained matrix with its \emph{rows permuted}
gives $1.000$, $0.984$ and $1.000$. With $A+1$ near-orthogonal vectors in
$d=64$, any scalar ranking of tokens is linearly realisable, so the probe
carries the order in its own weights. That result is withdrawn; the
hidden-state measurements, which use no probe, are unaffected.

\paragraph{The one direction that is causal is causal for \texttt{local}.}
Projecting the fitted order direction out of the embedding, against a
random-direction control at matched rank, costs an excess of
$+0.079$ to $+0.100$ nats on the overall loss ($4.3$--$6.2$ control standard
deviations) --- of which \texttt{local} absorbs $+0.279$ and \texttt{higher}
$+0.0035$.

\subsection{Statement}

Across twelve configurations spanning depth, width, input representation,
output factorisation, architecture family, objective and curriculum, the
level-3 probe never leaves the shuffled-label floor, while $1.87$ nats of
level-3 information are present in the prefix. We do not claim impossibility;
recurrent-depth architectures were tried only briefly, and the mechanistic
account below is a hypothesis.

\paragraph{Hypothesis (not established).}
Recovering the level-$\ell$ symbol requires marginalising over parses
consistent with the prefix. Inside--outside does this with an upward
\emph{and} a downward pass whose messages are coupled; a feedforward stack has
only the upward one and must approximate the composite in a fixed number of
layers. Consistently, depth from 2 to 8 buys $0.06$ nats and no change in
level-3 decodability.


\section{Part II: the constructed estimator and its verification}
\label{sec:hier-estimator}

\subsection{Construction}

No gradients are used. The rule tables $P_\ell[p,x,y]$ are estimated by:

\begin{enumerate}
\item \textbf{Spectral initialisation.} At level $\ell$, two observed child
pairs belong to the same parent if they occur in the same contexts. We build
the pair--context co-occurrence matrix, take its leading singular directions,
and $k$-means the pairs into $V$ groups. Where a level has fewer than four
nodes per sequence there is no horizontal context and we fall back to grouping
by shared child identity.

\item \textbf{Soft hand-off.} The posterior over parent symbols, not a hard
label, is propagated to the next level:
$q_\ell \propto \sum_{x,y} P_\ell[\cdot,x,y]\, q_{\ell-1}^{\mathrm{lo}}(x)\,
q_{\ell-1}^{\mathrm{hi}}(y)$.

\item \textbf{EM.} Inside--outside E-step, normalise-expected-counts M-step,
iterated; the sweep with the best held-out likelihood is kept.
\end{enumerate}

\subsection{Verification suite}

Each item below is an identity of the tree model class. Passing verifies the
implementation and that the fitted object is a genuine section; it does
\emph{not} test whether the data are tree-generated. Errors are maxima over
$256$ held-out sequences of Corpus~A.

\begin{center}
\begin{tabular}{llc}
\toprule
check & statement & max error \\
\midrule
V0 & EM never decreases the training log-likelihood & $+1.8\times10^{-2}$ (min step) \\
V1a & a free subtree has inside value exactly $1$ & $1.8\times10^{-15}$ \\
V1b & restriction $=$ marginalisation & $8.9\times10^{-16}$ \\
V1c & gluing $=$ independence join & $1.0\times10^{-14}$ \\
V2 & the hand-off is a sufficient statistic & $6.7\times10^{-16}$ \\
V4a & merge cost $= N\,I(\text{children};B \mid q(B))$ & $1.2\times10^{-14}$ \\
\bottomrule
\end{tabular}
\end{center}

\noindent
V1c is the concrete content of Theorem~\ref{thm:simpson-pullback} and
Construction~\ref{con:join}: with $\mu(u,v)$ the joint prior on the two top
children,
\begin{equation}
  P(x_{0:16}) \;=\; \sum_{u,v} \mu(u,v)\, \beta_L(u)\, \beta_R(v),
  \label{eq:glue}
\end{equation}
which we verify numerically rather than assume. Note $\mu$ is the \emph{joint},
not a product: the two halves are dependent through the root, and factorising
them breaks the identity.

V2 is the operational form of Definition~\ref{def:supports}: replacing each
completed subtree by its normalised inside vector and deleting its leaves
changes the predictive distribution by $<10^{-15}$.

V4a is Construction~\ref{con:merge}: merging symbols into their mixture row
costs exactly $\sum_s n_s\,\mathrm{KL}(R_s \Vert R_{\mathrm{class}})$, which is
the Bregman information of the merge and equals $N\,I(\text{children};B\mid
q(B))$ --- zero precisely when the merge is lossless.

\subsection{Estimator variants}

On Corpus~A, $2\times10^4$ training sequences:

\begin{center}
\begin{tabular}{lccc}
\toprule
variant & \texttt{local} & \texttt{higher} & held-out $\log p$/seq \\
\midrule
joint EM from random init      & 2.566 & 3.831 & $-51.18$ \\
bottom-up, frequency-rank clusters & 2.598 & 3.833 & $-51.45$ \\
spectral $+$ soft hand-off $+$ EM (10 sweeps) & 2.227 & 3.802 & $-48.23$ \\
spectral $+$ soft hand-off $+$ EM (50 sweeps) & 2.162 & 3.807 & $-47.75$ \\
stagewise $+$ vertical re-clustering & 2.240 & 3.802 & $-48.34$ \\
\midrule
transformer, $48$k steps        & 2.022 & 3.772 & --- \\
true grammar                    & 1.599 & 3.553 & $-41.21$ \\
\bottomrule
\end{tabular}
\end{center}

\noindent
Initialisation dominates. On Corpus~A the constructed estimator does not beat
the trained transformer; Section~\ref{sec:hier-regime} shows why, and
Section~\ref{sec:hier-splitmerge} closes the gap entirely on an identifiable
corpus.

\section{Part III: instruments for detecting hierarchy}
\label{sec:hier-instruments}

Three measurements, each with a calibration control and, where possible, a
planted positive.

\subsection{V6: where does the tree fail (token resolution)}

For each pair of leaf positions we compare empirical mutual information in the
data against the same statistic computed on datasets \emph{sampled from the
fitted model} at equal size (parametric bootstrap, $20$ replicates). Because
the same estimator is applied to both sides, plug-in bias cancels. Controls:
the true grammar against its own samples (calibration), and the true grammar
against the data (which also verifies that the exported rule tables are the
actual generator).

\paragraph{Planted positive.}
Copying leaf $7$ into leaf $8$ with probability $0.3$ creates a dependence
across the root that no tree of this shape carries. The detector reports
$z=+220$ at the planted pair against a largest other $z$ of $+6.5$; even the
true grammar is refuted there ($z=+390$). During training, the data-minus-model
gap at the planted pair stays at $+0.600$ across every sweep, while the
sibling gap closes from $+0.161$ to $+0.028$: training absorbs structure the
tree can hold and never touches structure it cannot.

\paragraph{Limitation.}
On Corpus~A, pairwise leaf MI across level-2 and root boundaries is $\approx
10^{-4}$ \emph{even under the true grammar}, while the plug-in bias floor at
$4\times10^4$ sequences is $\approx 0.055$. V6 therefore has no power above
level~1 there; the $0.3$ copy was detectable only because it was gross.

\subsection{V7: dependence between hierarchy units}

The fix is to test blocks, not tokens. Each block of $2^\ell$ leaves is decoded
to the level-$\ell$ symbol given by the argmax of \emph{that block's own}
inside vector, and mutual information is measured between adjacent decoded
blocks, with the same parametric bootstrap.

\paragraph{Disjointness invariant.}
Every statistic compares functions of disjoint leaf sets. Blocks share no
leaves, and decoding uses inside messages only: outside messages would import
the neighbour's evidence and make adjacent blocks appear dependent under any
model. Outside messages remain legitimate for \emph{fitting}; they are excluded
only from the test decoder.

\paragraph{Result (Corpus~A, $2\times10^4$ sequences).}
Decoded with the true grammar, blocks sharing a parent carry $\approx 1.0$ nat
at levels 1, 2 and 3, and blocks not sharing one carry $0.002$--$0.014$. The
fitted model captures $0.866/1.102 = 79\%$ at level 1, $0.014/0.994 = 1.4\%$ at
level 2 and $0.038/0.965 = 3.9\%$ at level 3. All calibration controls pass
($|z| \le 2.6$).

\subsection{$\beta_1$: is the block structure a tree at all}

Decode blocks as in V7 and build the graph on blocks whose pairwise MI exceeds
a threshold. Under a tree the dependency structure is nested: every threshold
cuts the graph into cliques (the blocks under one ancestor), and a clique
complex of cliques has no holes. Hence the prediction
\begin{equation}
  \text{tree} \;\Longrightarrow\; \beta_1 = 0 \ \text{at every threshold},
  \label{eq:betti}
\end{equation}
with $\beta_1 = \dim\ker\partial_1 - \operatorname{rank}\partial_2 =
(E-V+C) - \operatorname{rank}\partial_2$ computed exactly over $\mathrm{GF}(2)$
on the $2$-skeleton. A null band is obtained from datasets sampled from the
fitted tree.

\paragraph{Result and planted positive (Corpus~B).}
Unplanted: $\beta_1 = 0$ at all $25$ thresholds, null band $0$--$1$. Copying
block $0$ onto block $7$ and block $3$ onto block $4$, each with probability
$0.5$: $\beta_1 = 1$ at five thresholds with the null at $0$.

\paragraph{Two facts about the test.}
First, a \emph{single} extra coupling produces no cycle: it is a bridge between
two cliques, and $\beta_1$ counts independent loops, so at least two couplings
between distinct clique pairs are required. Second, planting must be done at
block level. Copying a single leaf creates ungrammatical pairs whose inside
vector is numerically zero, the decoded symbol degenerates, and the coupling
disappears; copying whole blocks keeps every production valid.


\section{Part IV: the regime determines whether anything is learnable}
\label{sec:hier-regime}

The failure of the estimator at levels 2--4 on Corpus~A is an identifiability
problem, not an algorithmic one. Level $\ell$ has $2^{L-\ell}$ nodes per
sequence and a table of $V\cdot c_\ell^2$ entries:

\begin{center}
\begin{tabular}{lcccc}
\toprule
 & level 1 & level 2 & level 3 & level 4 \\
\midrule
Corpus~A: observations per parameter & 1.4 & 0.31 & 0.15 & 0.08 \\
Corpus~B: observations per parameter & 1389 & 1563 & 781 & 391 \\
\bottomrule
\end{tabular}
\end{center}

\noindent
Corpus~B ($V=8$, $A=12$, $m=3$) was constructed so that every level is
estimable. On it the same estimator, with no algorithmic change, recovers the
hierarchy. Writing $c$ for the block dependence the fitted model's own symbols
capture and $g$ for the dependence its samples generate, and $z$ for the
bootstrap discrepancy:

\begin{center}
\begin{tabular}{lcccccc}
\toprule
 & \multicolumn{3}{c}{$30$ sweeps} & \multicolumn{3}{c}{$200$ sweeps} \\
\cmidrule(lr){2-4}\cmidrule(lr){5-7}
level & $c$ & $g$ & $z$ & $c$ & $g$ & $z$ \\
\midrule
1 & 0.4568 & 0.4574 & $-0.0$ & 0.4568 & 0.4530 & $+1.7$ \\
2 & 0.3582 & 0.3595 & $-0.2$ & 0.3804 & 0.3818 & $-0.2$ \\
3 & 0.0854 & 0.0507 & $+20.2$ & 0.1827 & 0.1715 & $+5.4$ \\
\bottomrule
\end{tabular}
\end{center}

\noindent
Level 3 more than doubles between $30$ and $200$ sweeps and its discrepancy
falls from $+20.2$ to $+5.4$; held-out log-likelihood improves from $-17.97$ to
$-17.51$ (floor $-15.32$) and was still descending at sweep $200$. The earlier
failure was therefore (i) an unestimable regime and (ii) undertraining, in that
order.

\paragraph{The symbols are not redundant.}
The merge sweep (V4) on Corpus~B: at levels 1--3 every merge costs real
conditional MI ($0.142$, $0.071$, $0.128$ nats/node at the first step) and
held-out likelihood falls monotonically. Level~4 is the exception: $8\to6$
costs $0.058$ nats/node with held-out change $+0.0000$, and $8\to2$ gives
$+0.0039$. The root alphabet is genuinely redundant; the others are not.

\paragraph{A defect, later resolved.}
At token resolution the fitted model still under-produced cross-boundary
dependence: $0.174$ against a true $0.282$ at the level-1 boundary, $0.0006$
against $0.0035$ at the root, unchanged from $30$ to $200$ sweeps and unchanged
when the training set was quadrupled from $5\times10^4$ to $2\times10^5$
sequences ($-17.5094$ against $-17.5113$ held-out). It is resolved in
Section~\ref{sec:hier-splitmerge}.

\section{Part V: what is not established}
\label{sec:hier-negative}

We record failed constructions, because several are attractive and the reasons
they fail are cheap to check.

\paragraph{The embedding probe (withdrawn).}
See Section~\ref{sec:hier-transformer}: a random matrix supports the same
probe accuracy.

\paragraph{Cohomological obstruction to gluing is vacuous here.}
For a single joint distribution on a fixed tree, every family of marginals
extends, by Theorem~\ref{thm:vorobev}. The \v{C}ech obstruction of
Abramsky--Mansfield--Barbosa~\cite{amb2011} is therefore identically zero. It
becomes meaningful only for genuinely overlapping contexts --- several parse
shapes over the same tokens, or models fitted on different corpora sharing
variables. The $\beta_1$ test of Section~\ref{sec:hier-instruments} is the
non-vacuous substitute.

\paragraph{Information cohomology computes entropy, not obstructions.}
Vigneaux~\cite{vigneaux2019}, following Baudot--Bennequin~\cite{bb2015}, shows
$H^1$ of the information site is generated by Shannon entropy. This
characterises which functionals exist; it does not measure failure to glue. A
proposed test --- verify the chain rule $H(X,Y)=H(X)+H(Y\mid X)$ at each node
--- cannot fail when all three terms are computed from the same empirical
joint: it is an algebraic identity.

\paragraph{The lattice quasi-morphism does not exist.}

\begin{proposition}\label{prop:padic}
Let $\rho:\mathbb{Q}\to L$ be order-preserving for the real order and constant
on each $p$-adic ball. Then $\rho$ is constant.
\end{proposition}

\begin{proof}
A $p$-adic ball $B(a,p^{-n}) \cap \mathbb{Q}$ contains $a + p^n m/k$ for all
$m \in \mathbb{Z}$ and all $k$ coprime to $p$, hence is dense in $\mathbb{R}$.
Let $B_1, B_2$ be disjoint balls with $\rho(B_1)=c_1$, $\rho(B_2)=c_2$. By
density choose $x<y<z$ with $x,z\in B_1$ and $y\in B_2$. Monotonicity gives
$c_1 \le c_2 \le c_1$, so $c_1=c_2$.
\end{proof}

\noindent
The underlying obstruction is Theorem~\ref{thm:ostrowski}: $\mathbb{R}$ and
$\mathbb{Q}_p$ are inequivalent completions of $\mathbb{Q}$, and the $p$-adic
norm of a probability measures divisibility, not magnitude ($|0.3|_3 = 1/3 <
|0.5|_3 = 1$). Real-valued rounding to the grid $p^{-n}\mathbb{Z}$ is a genuine
quasi-morphism with defect $\le p^{-n}$, but retains nothing $p$-adic.

\paragraph{Posterior entropy is the wrong confidence criterion.}
After stagewise fitting on Corpus~A the per-level posterior entropies are
$0.198$, $0.353$, $0.541$, $1.003$ nats against a maximum of $4.159$ --- all
low --- while V7 shows those levels capturing under $4\%$ of the available
dependence. The model is confident about labels that carry no structure. The
criterion that does separate them is $I(\text{parent};\text{children})$
computed from the rule table, which needs no ground truth and reads
$10.1, 8.1, 4.5, 5.3$ against true values $8.9, 9.0, 9.0, 9.0$.

\paragraph{Curriculum does not transfer.}
The curriculum result of Qiu et al.~\cite{qiu2024} (their ablation: $0.0\%$
baseline, $1.6\%$ with depth, $8.7\%$ with depth and reversal, $99.9\%$ with
sample weighting) does not reproduce here. A matched control with no easy
cases is marginally \emph{better} on both components
($2.043/3.789$ versus $2.085/3.803$), and level-3 decoding is unchanged.
Their failure mode was a broken easy component; ours is not.




\section{Closing the gap: the failure was basin selection}
\label{sec:hier-splitmerge}

\subsection{Diagnosis: start EM at the truth}

The shortfall of Section~\ref{sec:hier-regime} admits three explanations: the
model class is wrong, the data are insufficient, or the optimiser is in a bad
basin. The third is testable directly. Initialise EM at the true rule tables,
optionally blended with uniform at weight $\varepsilon$, and run it.

\begin{center}
\begin{tabular}{lcc}
\toprule
start & held-out $\log p$/seq & \texttt{local} / \texttt{higher} \\
\midrule
exact truth ($\varepsilon=0$)          & $-15.324$ & $0.3322 / 1.5833$ \\
truth $+\ 50\%$ uniform                & $-15.325$ & $0.3322 / 1.5834$ \\
truth $+\ 80\%$ uniform                & $-15.325$ & $0.3322 / 1.5834$ \\
truth $+\ 90\%$ uniform                & $-15.330$ & --- \\
truth $+\ 95\%$ uniform                & $-15.775$ & --- \\
truth $+\ 99\%$ uniform                & $-18.113$ & --- \\
uniform ($\varepsilon=1$)              & $-20.706$ & $0.5091 / 2.0791$ \\
\midrule
spectral init (ours)                   & $-17.511$ & $0.3952 / 1.7937$ \\
\bottomrule
\end{tabular}
\end{center}

\noindent
Three conclusions. The truth is a fixed point, so the model class is right and
the corpus is what we believe it to be. The basin is very wide: a start
carrying only $10\%$ of the truth, whose likelihood ($-31.2$) is worse than
uniform, climbs all the way back. And our spectral initialisation performs like
a $1\%$ admixture --- it fails to clear a low bar.

\subsection{Why the hard clustering cannot work}

The initialisation assigns each observed child pair to exactly one parent
cluster. Real grammars share pairs across parents: in Corpus~B, $4$ of the $16$
distinct level-1 pairs are generated by two different symbols ($2$ of $18$ at
level 2, $4$ of $16$ at levels 3 and 4). A hard partition therefore cannot
represent a quarter of the grammar however good the clustering is. Two
predictions follow, and both hold: sharpening the clustering makes matters
worse (using the sibling co-parent as context instead of linear adjacency gives
$-18.06$), and blurring it helps (blending with uniform at weight $0.9$ gives
$-15.87$).

\subsection{Split-merge: letting EM do the clustering}

The remaining gap is closed by moving the clustering \emph{inside} EM. For each
level in turn:

\begin{enumerate}
\item \textbf{Split}: duplicate every level-$\ell$ symbol with a multiplicative
perturbation, doubling that level's alphabet; the child axes of level $\ell+1$
double with it.
\item \textbf{EM}: re-fit the enlarged model.
\item \textbf{Merge}: greedily collapse back to $V$ symbols, cheapest first,
where the cost of a merge is the Bregman information of
Construction~\ref{con:merge}. The merge ranges over all $2V$ symbols, so it can
recombine across the split rather than undo it. This is the escape mechanism.
\item Accept the round only if held-out likelihood improves.
\end{enumerate}

\begin{center}
\begin{tabular}{lcc}
\toprule
stage & held-out $\log p$/seq & merge cost (nats/seq) \\
\midrule
warmup, $40$ EM sweeps        & $-17.187$ & --- \\
round 1, levels 1--4          & $-15.491$ & $0.160, 0.113, 0.008, 0.003$ \\
round 2, levels 1--4          & $-15.325$ & $0.042, 0.011, 0.007, 0.005$ \\
round 3, levels 1--4          & $-15.3253$ & $0.052, 0.020, 0.007, 0.004$ \\
\midrule
true grammar                  & $\mathbf{-15.3243}$ & --- \\
\bottomrule
\end{tabular}
\end{center}

\noindent
Total time $131$\,s. The final model is $0.001$ nats per sequence from the
truth, and \texttt{local}/\texttt{higher} read $0.3322/1.5835$ against the true
$0.3322/1.5833$. One round was rejected (round~2, level~4), so the accept test
is doing work.

\subsection{Verification of the recovered model}

Likelihood agreement is not by itself recovery of structure. All four
instruments confirm it independently.

\begin{center}
\begin{tabular}{lccc}
\toprule
 & before split-merge & after & true \\
\midrule
V3, sibling MI               & $1.5445$ & $1.5507$ & $1.5523$ \\
V3, level-1 boundary         & $0.1739$ & $0.2783$ & $0.2819$ \\
V3, level-2 boundary         & $0.0153$ & $0.0334$ & $0.0331$ \\
V3, root boundary            & $0.0006$ & $0.0034$ & $0.0035$ \\
V6, boundaries flagged       & $11$ of $18$ & $\mathbf{0}$ of $18$ & --- \\
V6, largest $|z|$            & $+58.6$ & $+1.7$ & --- \\
V7, level 3 ($c$ vs $g$)     & $0.1798/0.1720$ & $0.5863/0.5882$ & $0.6140$ \\
$\beta_1$ vs fitted null     & $6$ of $25$ & $\mathbf{0}$ of $25$ & --- \\
\bottomrule
\end{tabular}
\end{center}

\noindent
The $\beta_1$ result is the sharpest. Section~\ref{sec:hier-instruments}
reported $6$ of $25$ thresholds flagged against a null sampled from the fitted
model, all in the weak-edge window where the fit under-produced dependence,
while the same test against a null sampled from the \emph{true} grammar gave
$0$ of $25$ and still caught a planted violation at $9$ of $25$. The diagnosis
was therefore that $\beta_1$ was measuring the basin, not the data. With the
recovered model the flags vanish. The instrument fires on a bad fit, is silent
on a good one, and detects planted non-tree structure in both cases.

\subsection{The gap to the transformer}

On Corpus~A the exact floors are $\texttt{local}=1.599$, $\texttt{higher}=3.553$
and the best transformer of Section~\ref{sec:hier-transformer} reaches
$2.022/3.772$, with the level-3 latent symbol never decodable above the
shuffled-label floor. On Corpus~B, where every level is identifiable, the
gradient-free estimator reaches $0.3322/1.5835$ against floors of
$0.3322/1.5833$ --- that is, the floor to three decimal places --- in $131$
seconds without computing a single gradient. We emphasise that these are
different corpora and the comparison is therefore about \emph{kind}, not a
head-to-head benchmark: the transformer fails to represent the upper levels at
all under twelve configurations, whereas the estimator recovers them exactly
when the regime is identifiable, and its remaining failures are diagnosable
(basin selection) and curable (split-merge) rather than opaque.


\section{Why clustering is necessary at scale}
\label{sec:hier-scaling}

Sections~\ref{sec:hier-estimator}--\ref{sec:hier-regime} used $V=8$, where a
dense rule tensor is trivial. The motivation for cluster primes is what happens
at $V=10^3$.

\subsection{Cost of the dense recursion}

Let $N$ be the total number of leaf tokens in the corpus. A complete binary
tree over $N$ leaves has $N-1$ internal nodes \emph{in total}: the levels
partition the nodes, they do not replicate them. Each node evaluates
\begin{equation}
  T_\ell(b,n,p) \;=\; \sum_{x,y} P_\ell(p\to x,y)\,
  T_{\ell-1}(b,2n,x)\, T_{\ell-1}(b,2n{+}1,y),
  \label{eq:inside-cost}
\end{equation}
at a cost of $2V^3$ operations, so one inside pass costs $\approx 2NV^3$. With
outside and the M-step, a sweep is $\approx 4NV^3$. At $N=10^6$ and $V=1024$:
\begin{equation}
  C^{\text{dense}}_{\text{sweep}} \approx 4.3\times 10^{15},
  \qquad
  C^{\text{dense}}_{50\ \text{sweeps}} \approx 2.1\times 10^{17}.
\end{equation}
Memory splits into two parts that scale differently. The inside table is
$\approx 2NV$ floats, $8$~GB, which is large but not fatal; the rule tensors
are $V^3$ entries \emph{per level}, $4.3$~GB each and $69$~GB over $L=16$,
which is. Replacing $V$ by $K$ gives $4NK^3$ per sweep,
$\approx 8\times10^{11}$ over $50$ sweeps at $K=16$, with rule tensors of
$16$~KB per level, and a speedup of
\begin{equation}
  \left(\frac{V}{K}\right)^{3} = 64^3 = 262{,}144.
\end{equation}

\subsection{Cost ladder for $m$-ary grammars}

For an $m$-ary grammar the rule tensor has order $m+1$ and the per-node cost is
$V^{m+1}$ dense, $K^{m+1}$ compressed. Each extra tensor order multiplies the
compressed cost by $K$ and the dense cost by $V$, so the compression buys
$\log_K V$ orders of headroom --- at $V/K = 64$, about $2.5$ orders.

\begin{center}
\begin{tabular}{lccc}
\toprule
arity $m$ & tensor order & dense ($V=1024$) & compressed ($K=16$) \\
\midrule
2 (binary)  & 3 & $2.1\times10^{17}$ & $8\times10^{11}$ \\
3 (ternary) & 4 & $2.2\times10^{20}$ & $1.3\times10^{13}$ \\
4           & 5 & $2.2\times10^{23}$ & $2.1\times10^{14}$ \\
5           & 6 & $2.3\times10^{26}$ & $3.4\times10^{15}$ \\
\bottomrule
\end{tabular}
\end{center}

\subsection{Clustering belongs inside EM, not before it}

Section~\ref{sec:hier-splitmerge} revises what the initial clustering is for.
An externally imposed partition supplies a seed and nothing more: sharpening it
hurts, blurring it helps, and the symbols that finally matter are produced by
merge moves \emph{inside} EM, scored by the conditional mutual information of
Construction~\ref{con:merge} on expected counts from the E-step. The scaling
argument above is unaffected --- at $V=10^3$ the $V^3$ contraction must be
compressed --- but the compression should be discovered by split-merge rather
than fixed in advance.

\subsection{Compression is not free: the measured trade-off}

The scaling argument alone would suggest taking $K$ as small as convenient.
The merge sweep of Section~\ref{sec:hier-regime} measures what that costs. On
Corpus~B, per level, $I/\text{node}$ is the conditional mutual information paid
(Construction~\ref{con:merge}) and $\Delta$ is the change in held-out
log-likelihood per sequence:

\begin{center}
\begin{tabular}{lcccccccc}
\toprule
 & \multicolumn{2}{c}{level 1} & \multicolumn{2}{c}{level 2}
 & \multicolumn{2}{c}{level 3} & \multicolumn{2}{c}{level 4} \\
\cmidrule(lr){2-3}\cmidrule(lr){4-5}\cmidrule(lr){6-7}\cmidrule(lr){8-9}
$K$ & $I/\text{node}$ & $\Delta$ & $I/\text{node}$ & $\Delta$
    & $I/\text{node}$ & $\Delta$ & $I/\text{node}$ & $\Delta$ \\
\midrule
8 & 0.000 & $+0.000$ & 0.000 & $+0.000$ & 0.000 & $+0.000$ & 0.000 & $+0.000$ \\
6 & 0.142 & $+0.000$ & 0.071 & $-0.222$ & 0.128 & $-0.018$ & 0.058 & $+0.000$ \\
4 & 0.397 & $-0.970$ & 0.321 & $-0.716$ & 0.430 & $-0.121$ & 0.326 & $+0.003$ \\
3 & 0.626 & $-1.502$ & 0.559 & $-0.838$ & 0.663 & $-0.122$ & 0.489 & $+0.003$ \\
2 & 0.983 & $-1.957$ & 0.941 & $-0.880$ & 1.063 & $-0.124$ & 0.723 & $+0.004$ \\
\bottomrule
\end{tabular}
\end{center}

\noindent
Levels 1--3 pay for every merge in both currencies. Level~4 does not: $8\to2$
costs $0.72$ nats per node of conditional MI and \emph{improves} held-out
likelihood by $0.004$, because the root alphabet is redundant at this depth.
So the right $K$ is level-dependent and must be measured, not assumed; the
merge identity V4a supplies the measurement.

\subsection{What is projected and what is demonstrated}

The recovery results of Section~\ref{sec:hier-regime} are at $L=4$, $V=8$,
$2\times10^5$ sequences. The costs above are arithmetic, not measurements: we
have not run $L=16$, $V=1024$. One caveat is worth stating in advance. At
$L=16$ a sequence has $2^{16}=65{,}536$ leaves, so $N=10^6$ tokens is about
$15$ sequences, and the root level would be observed $15$ times. The
compression removes the compute barrier; it does not remove the
observations-per-parameter barrier of Section~\ref{sec:hier-regime}, which at
the top of a deep tree is governed by the number of \emph{sequences}, not
tokens.

\paragraph{A proposed argument we do not use.}
It is tempting to argue that attention is architecturally incapable of
representing a rank-$m$ rule tensor because each head is a rank-2 bilinear form
and CP rank $\Theta(V^{m-1})$ would require $V^{m-1}$ heads. We do not make
this claim. The network computes a predictive function, not the rule tensor;
the MLP sublayers, which our own readout shows carrying most of the order
construction ($+0.127$ at $\mathrm{mlp}_1$ against $+0.041$ at
$\mathrm{attn}_1$), are absent from the argument; and depth composes heads
nonlinearly. The transformer result of Section~\ref{sec:hier-transformer} is an
empirical negative over twelve configurations and does not need it.



\section{Open problems}
\label{sec:hier-open}

All of these concern the estimator or the instruments. None concerns
cohomology: Section~\ref{sec:hier-negative} records why the gluing obstruction
is vacuous for a single distribution on a fixed tree, and nothing since has
changed that.

\paragraph{O0. The cross-boundary defect, and why it is no longer open.}
Until Section~\ref{sec:hier-splitmerge} the binding problem was that the fitted
model under-produced leaf-pair dependence at the level-1 and level-2
boundaries ($0.174$ against a true $0.282$), and the natural reading was that
the many-to-one block decoder cannot force the induced leaf marginals to match.
That reading was wrong: the defect was an artefact of the basin, and it
disappears entirely once split-merge finds the right optimum. We record it
because the wrong reading was plausible and survived several rounds of
measurement. It remains untested on Corpus~A, where split-merge has not been
run.

\paragraph{O1. A regime-identifiability threshold.}
We measure that levels fail at $0.15$ observations per rule-table parameter and
succeed at $781$, but we have no threshold $\theta(V, L, \ell, m)$ separating
them. What is wanted is a statement of the form: with $n$ sequences of a
depth-$L$, $V$-symbol, $m$-rule grammar, the level-$\ell$ table is identifiable
up to relabelling with probability $1-\delta$ provided
$n \cdot 2^{L-\ell} \ge \theta$. The sample complexity of latent-tree models is
known in the spectral setting~\cite{choi2011}; the EM setting with our
initialisation is not covered.

\paragraph{O2. When does redundancy justify collapsing a level?}
Construction~\ref{con:merge} gives the exact cost of a merge, and
Section~\ref{sec:hier-regime} shows level~4 collapsing from $8$ to $2$ symbols
at a held-out \emph{gain}. What is missing is a rule: given the merge cost
$N\,I(\text{children};B\mid q(B))$ and the parameter count released, predict the
sign of the held-out change without measuring it. This is a model-selection
statement---the natural form is a comparison of the merge cost against the
effective number of parameters removed---and would turn the $K$-sweep from a
search into a calculation.

\paragraph{O3. Why rule-table mutual information beats posterior entropy.}
Section~\ref{sec:hier-negative} shows posterior entropy $H(q_\ell)$ is low
($0.20$--$1.00$ nats against a maximum of $4.16$) at levels that capture under
$4\%$ of the available dependence, while $I(\text{parent};\text{children})$
computed from the rule table separates them ($10.1, 8.1, 4.5, 5.3$ against true
$8.9, 9.0, 9.0, 9.0$). The observation is that a model can be certain about
labels that carry no structure. We have no theorem saying when the two
diverge, nor a threshold on $I(\text{parent};\text{children})$ that certifies a
level as determined --- which is what a stopping rule for stagewise fitting
needs.

\paragraph{O4. Power of the $\beta_1$ test.}
The detector finds two planted block couplings at $9$ of $25$ thresholds with
the null at $0$, and a single coupling never, because one extra edge is a
bridge rather than a loop. What is the detection threshold as a function of the
number of couplings $m$, their strength, the number of blocks and the sample
size? At minimum, $m \ge 2$ is necessary; the sufficient condition and the
power curve in coupling strength are unmeasured.

\paragraph{O5. A genuinely non-tree corpus.}
Every positive control here is a \emph{planted} coupling made by copying one
block onto another, which is destructive: it adds a dependence and removes the
one it overwrites. At level~2, with only four blocks, this makes a $4$-cycle
unreachable --- every plant we tried broke one of the four edges it needed. The
clean test is a corpus \emph{generated} from a non-tree structure, so that the
detector can be scored against a known ground truth rather than a perturbation.
This is a corpus-builder change, not a detector change.

\paragraph{O6. Split-merge at scale.}
The recovery result is at $L=4$, $V=8$, $2\times10^4$ sequences, in $131$
seconds. Section~\ref{sec:hier-scaling} gives the arithmetic for $L=16$,
$V=10^3$ but no measurement. The specific risk named there is that compression
removes the compute barrier and not the identifiability barrier: at $L=16$ a
sequence has $65{,}536$ leaves, so $10^6$ tokens is about $15$ sequences and the
root level is observed $15$ times. Whether split-merge tolerates that, and
whether the split/merge cycle stays stable when the alphabet is large enough
that the doubled model does not fit in memory, are untested.

\paragraph{O7. Recurrent-depth transformers.}
Section~\ref{sec:hier-transformer} covers twelve configurations, all with
distinct weights per layer. Weight sharing across depth --- a single block
applied $L$ times --- is the one architecture family in which iterated message
passing is representable at fixed parameter count, and it was tried only
briefly. A negative there would considerably strengthen the claim; a positive
would locate what the feedforward stack is missing.


\section{Inventory of theorems and constructions used}
\label{sec:hier-inventory}

\subsection*{Used and load-bearing}

\begin{theorem}[Inside--outside; Baker~\cite{baker1979}, Lari--Young~\cite{lari1990}]
\label{thm:inside}
For a PCFG on a fixed tree, the inside values $\beta$ and outside values
$\alpha$ computed by the two recursions give exact marginals: the posterior at
a node is $\propto \alpha\beta$, and the marginal likelihood is the root inside
value against the prior.
\end{theorem}

\begin{theorem}[EM ascent; Dempster--Laird--Rubin~\cite{dlr1977}]
\label{thm:em}
Each EM iteration does not decrease the observed-data log-likelihood. The proof
is Jensen's inequality on the $Q$-function; no geometric hypothesis is needed.
\end{theorem}

\begin{theorem}[Convergence; Wu~\cite{wu1983}]
\label{thm:wu}
Under regularity conditions the EM sequence converges to a stationary point of
the likelihood; the limit depends on the initialisation.
\end{theorem}

\begin{theorem}[Extension; Vorob'ev~\cite{vorobev1962}]
\label{thm:vorobev}
A family of consistent marginals on a hypergraph extends to a global
distribution for every consistent family if and only if the hypergraph is
acyclic. A fixed tree is acyclic.
\end{theorem}

\begin{theorem}[Independent pullback; Simpson~\cite{simpson2024}, Thm.~6.4 and Prop.~8.5]
\label{thm:simpson-pullback}
A presheaf on standard Borel probability spaces is an atomic sheaf if and only
if it maps independent squares to pullbacks; the independent pullback of a
cospan is the conditional product $\int P(\cdot\mid r{=}\omega)\otimes
P(\cdot\mid s{=}\omega)\,dP(\omega)$.
\end{theorem}

\begin{definition}[Supports; Simpson~\cite{simpson2024}, Def.~7.1]
\label{def:supports}
A support of a random variable is the terminal factorisation of it through a
sample space: the smallest randomness it depends on. $\mathrm{RV}(A)$ has
supports.
\end{definition}

\begin{construction}[Independence join; Wright~\cite{wright2023}, \S7.1]
\label{con:join}
For tables agreeing on a shared marginal, $(p\bowtie q)_{ijk} = p_{ij}q_{jk}/r_j$
is a weak pullback and is the maximum-entropy distribution with those
marginals. This is the same object as Theorem~\ref{thm:simpson-pullback} and is
the content of \eqref{eq:glue}.
\end{construction}

\begin{construction}[Merging as convex combination; Wright~\cite{wright2023}, \S7.2]
\label{con:merge-giry}
Conflicting categorical tables are merged by averaging distributions in the
Giry monad. The merged rule row is the mixture of the rows merged.
\end{construction}

\begin{theorem}[Bregman information; Bregman~\cite{bregman1967},
Banerjee et al.~\cite{banerjee2005}]
\label{thm:bregman}
For the Bregman divergence generated by negative entropy --- that is, for KL ---
the optimal representative of a cluster is its mean, and the clustering loss is
the Bregman information, which equals the mutual information between the data
and the cluster assignment.
\end{theorem}

\begin{construction}[Merge criterion]
\label{con:merge}
Merging latent symbols into their mixture row costs
$\sum_s n_s\,\mathrm{KL}(R_s\Vert R_{\mathrm{class}}) = N\,I(\text{children};B
\mid q(B))$, zero exactly when the merge is lossless. Losslessness is strong
lumpability in the sense of Kemeny--Snell~\cite{kemeny1960}: the merged rows
must be identical.
\end{construction}

\begin{construction}[Agglomerative merging; Slonim--Tishby~\cite{slonim1999}]
\label{con:aib}
Greedy merging of the pair with least information loss is the agglomerative
information bottleneck; the split-merge PCFG estimator of Petrov et
al.~\cite{petrov2006} is the grammar-specific analogue.
\end{construction}

\begin{theorem}[Ostrowski~\cite{ostrowski1916}]
\label{thm:ostrowski}
Every nontrivial absolute value on $\mathbb{Q}$ is equivalent either to the
real absolute value or to a $p$-adic one, and the resulting completions
$\mathbb{R}, \mathbb{Q}_2, \mathbb{Q}_3, \dots$ are pairwise inequivalent.
\end{theorem}

\begin{construction}[Bias correction]
\label{con:mm}
Plug-in mutual information on a $K\times L$ table from $N$ samples is biased
upward by $\approx (K-1)(L-1)/2N$ (Miller--Madow~\cite{miller1955}). The
tracker applies the correction; the bootstrap does not need it, because the
same estimator is applied to data and to model samples.
\end{construction}

\begin{construction}[$\beta_1$ of the clique complex]
\label{con:betti}
$\beta_1 = \dim\ker\partial_1 - \operatorname{rank}\partial_2$, with
$\dim\ker\partial_1 = E - V + C$ the cycle rank, computed over
$\mathrm{GF}(2)$. See \eqref{eq:betti}.
\end{construction}

\subsection*{Consulted, with the reason they are not used}

\begin{itemize}
\item \textbf{Abramsky--Brandenburger~\cite{ab2011}, Abramsky--Mansfield--Barbosa~\cite{amb2011}}:
      sheaf-theoretic contextuality and its \v{C}ech obstruction. Vacuous for a
      single distribution on one fixed tree
      (Theorem~\ref{thm:vorobev}); relevant only for overlapping contexts.
\item \textbf{Vigneaux~\cite{vigneaux2019}, Baudot--Bennequin~\cite{bb2015}}:
      information cohomology. $H^1$ is generated by entropy; this is a
      characterisation, not an obstruction theory.
\item \textbf{Four-point condition / latent-tree distances
      (Buneman~\cite{buneman1971}; Choi et al.~\cite{choi2011})}: the additive
      information distance assumes equal state counts and full-rank joints.
      With $A=48$ leaves and $V=64$ latents neither holds, so additivity is not
      guaranteed and the test would not be trustworthy. The parametric
      bootstrap is valid regardless of state counts.
\item \textbf{Chow--Liu~\cite{chowliu1968}}: gives the best tree over
      \emph{observed} variables; our structure is latent.
\item \textbf{Amari~\cite{amari2016}}: $e$- and $m$-projections. The E-step is
      the $e$-projection and the M-step the $m$-projection; EM monotonicity
      follows from Theorem~\ref{thm:em} and does not require the Pythagorean
      identity.
\item \textbf{Tucker~\cite{tucker1966} / non-negative Tucker}: replacing
      $P_\ell$ by a rank-$R$ core reduces $V^3$ parameters to $R^3+3VR$. Not
      needed at $V=8$; becomes necessary at $V=10^3$, where a dense table is
      $10^9$ entries. Proposed, not implemented.
\item \textbf{Ili\'c-Stepi\'c et al.~\cite{ilic2014}}: conditional $p$-adic
      probability logic. Establishes that Monna--Springer boundedness needs the
      range restricted to $\mathbb{Z}_p$ or the propositional alphabet kept
      finite, and states explicitly that $\mathbb{Q}_p$ admits no compatible
      order --- which is why Assumption~3 (monotone along the tree, a statement
      about the real order) and $p$-adic valuation cannot be imposed together.
\item \textbf{Gauthier~\cite{gauthier2014}}: recasts probability spaces as
      Grothendieck sites and then removes time and probability, so the absence
      of uncertainty is definitional rather than derived. No algorithm; nothing
      bearing on estimation.
\item \textbf{Rute~\cite{rute2019}}: computable measure theory and algorithmic
      randomness. Relevant only if a \emph{constructive} convergence statement
      for EM were wanted in place of Theorem~\ref{thm:wu}.
\item \textbf{Ostrowski~\cite{ostrowski1916}, Monna--Springer~\cite{ms1963}}:
      used only to establish Proposition~\ref{prop:padic} and the
      incompatibility of $p$-adic valuation with the real ordering of
      probabilities.
\end{itemize}


\section{Reproduction: commands and the result each produced}
\label{sec:hier-repro}

Every number in this section came from the commands below, in this order.
Corpus~A is the inherited configuration ($V=64$, $A=48$); Corpus~B is the
identifiable one ($V=8$, $A=12$).

\subsection*{Corpora}

\begin{verbatim}
python3 build_corpus_syn.py --out /tmp --synonym --reverse-valid \
        --nleaf 48 --syn-frac 0.35 --nsym 64 --nrules 12 \
        --poset-density 0.25 --n-train 200000
\end{verbatim}
Corpus~A. Observations per rule-table parameter: $1.4, 0.31, 0.15, 0.08$ at
levels 1--4. Held-out bigram $3.534$.

\begin{verbatim}
python3 build_corpus_syn.py --out /tmp/cH --depth 4 --nsym 8 --nrules 3 \
        --nleaf 12 --reverse-valid --n-train 200000 --n-val 5000
\end{verbatim}
Corpus~B. Observations per parameter: $1389, 1563, 781, 391$. Exact entropy
$1.160$ nats/token.

\subsection*{Exact floors and latent-parse entropies}

\begin{verbatim}
python3 floor_exact.py --data /tmp --samples 500
\end{verbatim}
Corpus~A: $H_{\mathrm{local}}=1.6034$, $H_{\mathrm{higher}}=3.5527$,
$H(x_0)=3.8161$. Latent-symbol entropy given the prefix, averaged over higher
positions: $2.784, 2.292, 2.290, 2.791$ against chance $4.159$ --- so $1.87$
nats of level-3 information are recoverable.

\subsection*{Transformer arms (Corpus~A, $48{,}000$ steps each)}

\begin{verbatim}
python3 rhm16.py --mlp --layers 2 --freeze-emb --steps 48000 --seed 1 --save L2.pt
python3 rhm16.py --mlp --layers 4 --freeze-emb --steps 48000 --seed 1 --save L4.pt
python3 rhm16.py --mlp --layers 8 --freeze-emb --steps 48000 --seed 1 --save L8.pt
python3 rhm16.py --mlp --layers 4 --d-model 128 --freeze-emb --steps 48000 \
        --seed 1 --save W128.pt
python3 rhm17.py --mlp --layers 8 --freeze-emb --steps 48000 --seed 1 --save A0.pt
python3 rhm17.py --mlp --layers 8 --freeze-emb --steps 48000 --seed 1 \
        --fresh-sample --save A1.pt
python3 rhm17.py --mlp --layers 8 --freeze-emb --steps 48000 --seed 1 \
        --aux-levels 3,4 --save A2.pt
python3 tree_cnn2.py --steps 48000 --seed 1 --save C2.pt
python3 curriculum.py --steps 48000 --seed 1 --save K1.pt
python3 curriculum.py --steps 48000 --seed 1 --simple-frac 0 --save K0.pt
python3 latent_head.py --steps 48000 --seed 1 --save M1.pt
python3 looped.py --loops 4 --steps 48000 --seed 1 --save P4.pt
\end{verbatim}
\texttt{local}/\texttt{higher} at step $48{,}000$, floors $1.603/3.553$:
$2.173/3.831$ (2 layers), $2.078/3.797$ (4), $2.037/3.755$ (8 frozen),
$2.022/3.772$ (8 frozen, fresh sampling), $2.039/3.770$ ($d{=}128$),
$2.088/3.777$ (auxiliary supervision), $2.104/3.827$ (dilated CNN),
$2.085/3.803$ and $2.043/3.789$ (curriculum and its control),
$3.227/3.835$ (latent head, stopped at $34$k).

\subsection*{Probes and controls (Corpus~A)}

\begin{verbatim}
python3 parse_state.py --load A1.pt --layers 8
python3 parse_state.py --load A2.pt --layers 8
python3 null_probe.py --load s1.pt
python3 ablate_w.py --load s1.pt
\end{verbatim}
Level-$\ell$ decoding from $h[t]$: $0.479/0.213$ (level 1),
$0.072/0.064$ (2), $0.028/0.024$ (3), $0.020/0.018$ (4), against shuffled
$0.021$--$0.034$ and chance $0.0156$. The auxiliary head trained on
ground-truth labels reaches $0.056$ (level 3) and $0.019$ (level 4).
\texttt{null\_probe} returns $1.000$ on a random matrix, $0.984$ on an
orthonormal one and $1.000$ on the trained matrix with rows permuted: the
embedding probe is retired. \texttt{ablate\_w} gives an excess of
$+0.079$--$+0.100$ nats over a matched random-direction control
($4.3$--$6.2$\,sd), of which \texttt{local} absorbs $+0.279$ and
\texttt{higher} $+0.0035$.

\subsection*{Estimator and verification suite}

\begin{verbatim}
python3 estimator2.py --data /tmp --n-train 20000 --polish 50

python3 estimator4.py --data /tmp --n-train 20000 --polish 50 --skip-v4 \
        --boot 20 --boot-n 40000 --save fit20k.npz

python3 estimator4.py --data /tmp/cH --n-train 200000 --polish 200 \
        --boot 20 --boot-n 40000 --skip-v4 --save /tmp/cH/fit200k.npz
\end{verbatim}
Corpus~A: $2.162/3.807$, held-out $-47.75$ (the numpy port of
\texttt{estimator4} reproduces \texttt{estimator2} to $10^{-3}$).
Corpus~B at $2\times10^5$ sequences: $0.3952/1.7937$, held-out $-17.511$ ---
identical to the $5\times10^4$ run ($-17.509$), so the residual gap is not
sample-limited. Verification: V1a $1.8\times10^{-15}$, V1b $8.9\times10^{-16}$,
V1c $1.0\times10^{-14}$, V2 $6.7\times10^{-16}$, V4a $1.2\times10^{-14}$.

\begin{verbatim}
python3 estimator4.py --data /tmp/cH --n-train 50000 --polish 30 \
        --boot 4 --boot-n 40000 --ks 8,6,4,3,2
\end{verbatim}
Merge sweep: levels 1--3 pay in both currencies (first-merge $I/\text{node}$
$0.142, 0.071, 0.128$; held-out falls monotonically to $-1.96, -0.88, -0.12$ at
$K=2$). Level~4 is redundant: $8\to2$ costs $0.72$ nats/node and \emph{improves}
held-out by $+0.004$.

\subsection*{Diagnosis of the local optimum}

\begin{verbatim}
python3 estimator4.py --data /tmp/cH --n-train 20000 --polish 15 --init-true 0.0 \
        --boot 3 --boot-n 4000 --skip-v4
python3 estimator4.py --data /tmp/cH --n-train 20000 --polish 40 --init-true 0.9 \
        --boot 3 --boot-n 4000 --skip-v4
python3 estimator4.py --data /tmp/cH --n-train 20000 --polish 60 --anneal 5 \
        --boot 3 --boot-n 4000 --skip-v4
python3 estimator4.py --data /tmp/cH --n-train 20000 --polish 60 --sib-ctx \
        --boot 3 --boot-n 4000 --skip-v4
python3 estimator4.py --data /tmp/cH --n-train 20000 --polish 300 --init-mix 0.9 \
        --boot 3 --boot-n 4000 --skip-v4
\end{verbatim}
Truth is a fixed point ($-15.324$, unchanged over $15$ sweeps). Truth blended
with $90\%$ uniform recovers ($-15.330$); with $99\%$ it does not ($-18.113$).
Annealing at $T_0\in\{2,5,20\}$ all give $-17.632$. Sibling context is
\emph{worse} ($-18.06$). Blending the init with uniform at $0.9$ gives
$-15.865$.

\subsection*{Split-merge}

\begin{verbatim}
python3 splitmerge.py --data /tmp/cH --n-train 20000 --warmup 40 --rounds 3 \
        --split-sweeps 20 --post-sweeps 20 --save /tmp/cH/sm.npz

python3 estimator4.py --data /tmp/cH --load-P /tmp/cH/sm.npz \
        --n-train 20000 --boot 20 --boot-n 40000 --skip-v4
\end{verbatim}
Held-out $-15.3253$ against the true grammar's $-15.3243$, in $131$\,s.
\texttt{local}/\texttt{higher} $=0.3322/1.5835$ against $0.3322/1.5833$.
Verification of the recovered model: V3 matches the true column at every
boundary, V6 flags $0$ of $18$ with largest $|z|=1.7$, V7 tracks at all three
levels with $|z|\le1.6$.

\subsection*{Obstruction and $\beta_1$ detectors}

\begin{verbatim}
python3 estimator4.py --data /tmp/cH --n-train 50000 --polish 30 \
        --plant 7,8,0.3 --boot 20 --boot-n 40000 --skip-v4

python3 betti.py --data /tmp/cH --load-P /tmp/cH/fit.npz --level 1 \
        --n 40000 --boot 20 --null true --show-mi
python3 betti.py --data /tmp/cH --load-P /tmp/cH/fit.npz --level 1 \
        --n 40000 --boot 20 --null true --plant-block "0,7,0.5;3,4,0.5"
python3 betti.py --data /tmp/cH --load-P /tmp/cH/fit.npz --level 1 \
        --n 40000 --boot 20 --null fitted
python3 betti.py --data /tmp/cH --load-P /tmp/cH/sm.npz --level 1 \
        --n 40000 --boot 20 --null fitted
\end{verbatim}
Planted leaf copy: $z=+220$ at the planted boundary against $+6.5$ elsewhere;
the true grammar is refuted there at $z=+390$. $\beta_1$ against a true-grammar
null: $0$ of $25$ unplanted, $9$ of $25$ with two planted block couplings.
Against a fitted null: $6$ of $25$ for the local-optimum fit, $0$ of $25$ for
the split-merge fit.

\paragraph{Note on settings.} \texttt{--boot} below about $10$ makes the
bootstrap standard deviation unreliable and the control checks will fail
spuriously; \texttt{--boot-n} below about $4\times10^4$ makes the Miller--Madow
correction dominate the $64\times64$ block statistics. The small values above
appear only in diagnostic runs where V6/V7 are not being read.

\begin{thebibliography}{99}

\bibitem{ab2011} S.~Abramsky and A.~Brandenburger.
  The sheaf-theoretic structure of non-locality and contextuality.
  \emph{New J. Phys.} 13:113036, 2011.

\bibitem{amb2011} S.~Abramsky, S.~Mansfield, and R.~S.~Barbosa.
  The cohomology of non-locality and contextuality.
  \emph{EPTCS} 95:1--14, 2012.

\bibitem{amari2016} S.~Amari. \emph{Information Geometry and Its
  Applications}. Springer, 2016.

\bibitem{baker1979} J.~K.~Baker. Trainable grammars for speech recognition.
  \emph{Speech Communication Papers, 97th Meeting of the ASA}, 1979.

\bibitem{banerjee2005} A.~Banerjee, S.~Merugu, I.~S.~Dhillon, and J.~Ghosh.
  Clustering with Bregman divergences. \emph{JMLR} 6:1705--1749, 2005.

\bibitem{bb2015} P.~Baudot and D.~Bennequin. The homological nature of entropy.
  \emph{Entropy} 17(5):3253--3318, 2015.

\bibitem{bregman1967} L.~M.~Bregman. The relaxation method of finding the
  common point of convex sets. \emph{USSR Comput. Math. Math. Phys.}
  7(3):200--217, 1967.

\bibitem{buneman1971} P.~Buneman. The recovery of trees from measures of
  dissimilarity. In \emph{Mathematics in the Archaeological and Historical
  Sciences}, 387--395, 1971.

\bibitem{cagnetta2024} F.~Cagnetta, L.~Petrini, U.~M.~Tomasini, A.~Favero, and
  M.~Wyart. How deep neural networks learn compositional data: the Random
  Hierarchy Model. \emph{Phys. Rev. X} 14:031001, 2024.

\bibitem{choi2011} M.~J.~Choi, V.~Y.~F.~Tan, A.~Anandkumar, and A.~S.~Willsky.
  Learning latent tree graphical models. \emph{JMLR} 12:1771--1812, 2011.

\bibitem{chowliu1968} C.~K.~Chow and C.~N.~Liu. Approximating discrete
  probability distributions with dependence trees. \emph{IEEE Trans. Inform.
  Theory} 14(3):462--467, 1968.

\bibitem{dlr1977} A.~P.~Dempster, N.~M.~Laird, and D.~B.~Rubin. Maximum
  likelihood from incomplete data via the EM algorithm. \emph{JRSS B}
  39(1):1--38, 1977.

\bibitem{gauthier2014} R. Gauthier. Algebraic stochastic calculus.
  \texttt{arXiv:1407.6784}, 2014.

\bibitem{ilic2014} A. Ili\'c-Stepi\'c, Z. Ognjanovi\'c, and N. Ikodinovi\'c.
  Conditional $p$-adic probability logic. \emph{Int. J. Approx. Reasoning}
  55:1843--1865, 2014.

\bibitem{kemeny1960} J.~G.~Kemeny and J.~L.~Snell. \emph{Finite Markov Chains}.
  Van Nostrand, 1960.

\bibitem{lari1990} K.~Lari and S.~J.~Young. The estimation of stochastic
  context-free grammars using the inside--outside algorithm.
  \emph{Computer Speech and Language} 4(1):35--56, 1990.

\bibitem{miller1955} G.~A.~Miller. Note on the bias of information estimates.
  In \emph{Information Theory in Psychology}, 95--100, 1955.

\bibitem{ms1963} A.~Monna and T.~Springer. Int\'egration non-archim\'edienne.
  \emph{Indag. Math.} 25:634--653, 1963.

\bibitem{ostrowski1916} A.~Ostrowski. \"Uber einige L\"osungen der
  Funktionalgleichung $\varphi(x)\varphi(y)=\varphi(xy)$.
  \emph{Acta Math.} 41:271--284, 1916.

\bibitem{petrov2006} S.~Petrov, L.~Barrett, R.~Thibaux, and D.~Klein. Learning
  accurate, compact, and interpretable tree annotation. \emph{ACL}, 2006.

\bibitem{qiu2024} X.~Qiu et al. Dissecting multiplication in transformers.
  \texttt{arXiv:2407.15360}, 2024.

\bibitem{rute2019} J. Rute. On the close interaction between algorithmic
  randomness and constructive/computable measure theory.
  \texttt{arXiv:1812.03375}, 2019.

\bibitem{simpson2024} A.~Simpson. Equivalence and conditional independence in
  atomic sheaf logic. \texttt{arXiv:2405.11073}, 2024.

\bibitem{slonim1999} N.~Slonim and N.~Tishby. Agglomerative information
  bottleneck. \emph{NIPS}, 1999.

\bibitem{tucker1966} L.~R.~Tucker. Some mathematical notes on three-mode factor
  analysis. \emph{Psychometrika} 31:279--311, 1966.

\bibitem{vigneaux2019} J.~P.~Vigneaux. \emph{Topology of Statistical Systems: A
  Cohomological Approach to Information Theory}. PhD thesis, Universit\'e Paris
  Diderot, 2019.

\bibitem{vorobev1962} N.~N.~Vorob'ev. Consistent families of measures and their
  extensions. \emph{Theory Probab. Appl.} 7(2):147--163, 1962.

\bibitem{wright2023} G.~Wright. An algebraic perspective on computing with
  data. PhD thesis, 2023.

\bibitem{wu1983} C.~F.~J.~Wu. On the convergence properties of the EM
  algorithm. \emph{Ann. Statist.} 11(1):95--103, 1983.

\end{thebibliography}

\end{document}
